% 附录 E — 实用程序（精译重写）
\biappendix{实用程序}{Utility Programs}
\label{chap:util}

\section{引言}
\subsection*{Introduction}

本附录介绍 RedHawk 安装目录 \texttt{bin} 下的关键命令行实用程序。

\begin{itemize}
\item \textbf{vcdtrans} — 从 VCD 文件统计各网络翻转次数，生成 toggle 文件供功耗与动态分析使用。
\item \textbf{vcdscan} — 从 VCD 计算每周期峰值/平均功耗，确定最坏功耗时间窗口。
\item \textbf{fsdbtrans} — 从 FSDB 文件生成各网络 toggle 文件（功能同 vcdtrans）。
\item \textbf{ircx2tech} — 将 iRCX 工艺文件转换为 RedHawk .tech 工艺文件。
\item \textbf{rhtech} — 从代工厂 nxtgrd/ITF 等标准工艺文件生成 RedHawk .tech 文件。
\item \textbf{gds2rh/gds2def} — 从 GDSII 数据生成 DEF/LEF，用于 RedHawk 数据库集成。
\item \textbf{pt2timing} — PrimeTime TCL 接口：生成各实例的转换时间与 timing window。
\item \textbf{sim2iprof} — 从仿真输出生成电流轮廓（read/write/standby 波形）。
\item \textbf{aplreader} — 将电流轮廓数据转换为标准 APL 输出格式。
\item \textbf{aplcdev2pwc} — 校验 cdev 数据有效性并转换为 pwcdev 格式。
\item \textbf{aplcopy} — 复制并编辑 *spiprof、*cdev、*pwcdev 文件中的数据。
\item \textbf{aplchk} — 对 APL 电容与 PWC 数据进行有效性检查。
\item \textbf{clampviewer} — Linux 可执行工具，预览 ESD clamp 的 I-V 曲线数据。
\end{itemize}

\section{vcdtrans}
\subsection*{vcdtrans}

从 VCD 文件统计各网络翻转次数，生成 toggle 文件供功耗与动态分析使用。

GSR 关键字 ：用于设置与  相关的 RedHawk 运行参数。

GSR 关键字 ：用于设置与  相关的 RedHawk 运行参数。

\begin{lstlisting}
            file and the outputs are individual block toggle files called <block>.toggle. The
\end{lstlisting}

GSR 关键字 ：用于设置与  相关的 RedHawk 运行参数。

GSR 关键字 ：用于设置与  相关的 RedHawk 运行参数。

\begin{lstlisting}
                vcdtrans <input_file_name> -o <output_file_name> -c
                    -s <start_time_in_ps> -e <end_time_in_ps>
                    -w “<logical_module_name>” “<physical_module_name>”
\end{lstlisting}

GSR 关键字 ：用于设置与  相关的 RedHawk 运行参数。

\begin{lstlisting}
                    <input_file_name> : specifies the VCD file.
                    -o <output_file_name> : specifies the toggle file.
                    -c : enables case-insensitive checking.
                    -s <start_time_in_ps> : specifies the start time, an integer with a default of t=0.
                    -e <end_time_in_ps> : specifies the end time, an integer with default of full
\end{lstlisting}

仿真时间

\begin{lstlisting}
                    -w “<logical_module_name>” : defines the logical module name from VCD
\end{lstlisting}

\section{vcdscan}
\subsection*{vcdscan}

从 VCD 计算每周期峰值/平均功耗，确定最坏功耗时间窗口。

GSR 关键字 ：用于设置与  相关的 RedHawk 运行参数。

\begin{lstlisting}
                     “<physical_module_name>” : defines the physical module name from DEF.
\end{lstlisting}

GSR 关键字 ：用于设置与  相关的 RedHawk 运行参数。

GSR 关键字 ：用于设置与  相关的 RedHawk 运行参数。

GSR 关键字 ：用于设置与  相关的 RedHawk 运行参数。（默认：located in the adsPower directory）

GSR 关键字 ：用于设置与  相关的 RedHawk 运行参数。

\begin{lstlisting}
                   vcdscan <VCD_file_name> -f <cycle_time_in_ps>
                        -w “<logical_module_name>” “<physical_module_name>”
                        -msg <msg_file_name> -o <output_file_name>
                        -d <input_file_dir_path> -cmd <command_file>
                        -s <start_time> -e <end_time>
                        -a <sim_frame_start_time> <sim_frame_end_time> -tt
\end{lstlisting}

GSR 关键字 ：用于设置与  相关的 RedHawk 运行参数。

\begin{lstlisting}
                     <VCD_file_name> : specifies the VCD file from Verilog simulation.
                     -f <cycle_time_in_ps> : specifies the frame size or cycle time in picoseconds.
                     -s <start_time> -e <end_time> : defines the start and end times for a power
\end{lstlisting}

GSR 关键字 ：用于设置与  相关的 RedHawk 运行参数。

\begin{lstlisting}
                     -a <sim_frame_start_time> <sim_frame_end_time> : defines the start and end
\end{lstlisting}

GSR 关键字 ：用于设置与  相关的 RedHawk 运行参数。

\begin{lstlisting}
                     -w “<logical_module_name>” : specifies the logical module name from VCD
                     “<physical_module_name>” : specifies the physical module name from DEF.
                     -msg <message_file_name> : defines where the messages are saved. The
\end{lstlisting}

GSR 关键字 ：用于设置与  相关的 RedHawk 运行参数。（默认：is print to screen）

\begin{lstlisting}
                     -o <output_file_name> : defines where the output file containing frame-by-frame
                         power data are saved. The default is <design>.pwr.
                     -d <input_file_dir_path> : defines the relative or absolute path to the input files.
\end{lstlisting}

GSR 关键字 ：用于设置与  相关的 RedHawk 运行参数。（默认：is the adsPower dir）

\begin{lstlisting}
                     -cmd <command file> : defines a command file containing the vcdscan command
\end{lstlisting}

GSR 关键字 ：用于设置与  相关的 RedHawk 运行参数。

\begin{lstlisting}
                     -tt : selects use of true timing results from vcd file for simulation
\end{lstlisting}

GSR 关键字 ：用于设置与  相关的 RedHawk 运行参数。

\begin{lstlisting}
                     -w "test/top" "" -o design.pwr
\end{lstlisting}

GSR 关键字 ：用于设置与  相关的 RedHawk 运行参数。

\begin{lstlisting}
             NOTE:            Remember the peak power times (<FROM>, <TO>) and the worst-case
\end{lstlisting}

GSR 关键字 ：用于设置与  相关的 RedHawk 运行参数。

GSR 关键字 ：用于设置与  相关的 RedHawk 运行参数。

\begin{lstlisting}
            Dynamic > VCD-based Dynamic IR-drop Analysis command. The command prompts
            for simulation parameters, such as the peak power periods (<FROM>, <TO>) and the
\end{lstlisting}

GSR 关键字 ：用于设置与  相关的 RedHawk 运行参数。

\section{fsdbtrans}
\subsection*{fsdbtrans}

从 FSDB 文件生成各网络 toggle 文件（功能同 vcdtrans）。

GSR 关键字 ：用于设置与  相关的 RedHawk 运行参数。

GSR 关键字 ：用于设置与  相关的 RedHawk 运行参数。

\begin{lstlisting}
            the vcd utilities and proceed directly to running Dynamic > VCD-based Dynamic IR-
\end{lstlisting}

GSR 关键字 ：用于设置与  相关的 RedHawk 运行参数。

GSR 关键字 ：用于设置与  相关的 RedHawk 运行参数。

GSR 关键字 ：用于设置与  相关的 RedHawk 运行参数。

\begin{lstlisting}
            file and the outputs are individual block toggle files with names <block>.toggle. The
\end{lstlisting}

GSR 关键字 ：用于设置与  相关的 RedHawk 运行参数。

GSR 关键字 ：用于设置与  相关的 RedHawk 运行参数。

\begin{lstlisting}
                fsdbtrans <input_file_name> -o <output_file_name> -c
                    -s <start_time_in_ps> -e <end_time_in_ps>
                    -w “<logical_module_name>” “<physical_module_name>”
\end{lstlisting}

GSR 关键字 ：用于设置与  相关的 RedHawk 运行参数。

\begin{lstlisting}
                    <input_file_name> : specifies the FSDB file.
                    -o <output_file_name> : specifies the toggle file.
                    -c : enables case-insensitive checking.
                    -s <start_time> : defines the start of scan time sequence in ps
                    -e <end_time> : defines the end of scan time sequence in ps
                    -w “<logical_module_name>” : specifies the logical module name from FSDB
                    “<physical_module_name>” : specifies the physical module name from DEF.
\end{lstlisting}

GSR 关键字 ：用于设置与  相关的 RedHawk 运行参数。

\section{ircx2tech}
\subsection*{ircx2tech}

将 iRCX 工艺文件转换为 RedHawk .tech 工艺文件。

GSR 关键字 ：用于设置与  相关的 RedHawk 运行参数。

\begin{lstlisting}
             ircx2tech <options>
\end{lstlisting}

GSR 关键字 ：用于设置与  相关的 RedHawk 运行参数。

\begin{lstlisting}
                  ircx2tech -i <input_file> -o <output_file> -v [min|typ|max]
                      ?-m <layer_mapping_file>? ?-g captab ?
                      ? -pe <em_rule_file>? ? -ct <connectivity file> ?
                      ? -rcm <RC_mapping file> ? ? -e <input_ircx_file_EM> ?
\end{lstlisting}

GSR 关键字 ：用于设置与  相关的 RedHawk 运行参数。

\begin{lstlisting}
                    -i: specifies the input iRCX file
                    -o: specifies the name for the output tech file
                    -v: specifies a via resistance table with a different RC corner iRCX file
\end{lstlisting}

GSR 关键字 ：用于设置与  相关的 RedHawk 运行参数。

\begin{lstlisting}
                    -m: specifies the name of the layer mapping file that contains the layer mapping
\end{lstlisting}

GSR 关键字 ：用于设置与  相关的 RedHawk 运行参数。

\begin{lstlisting}
                    -g captab: if specified, RedHawk calls $APACHEROOT/bin/captab to generate the
\end{lstlisting}

GSR 关键字 ：用于设置与  相关的 RedHawk 运行参数。

\begin{lstlisting}
                    -pe : specifies the em rule file which contains the EM rule per layer.The rule
\end{lstlisting}

GSR 关键字 ：用于设置与  相关的 RedHawk 运行参数。

\section{rhtech}
\subsection*{rhtech}

从代工厂 nxtgrd/ITF 等标准工艺文件生成 RedHawk .tech 文件。

GSR 关键字 ：用于设置与  相关的 RedHawk 运行参数。

GSR 关键字 ：用于设置与  相关的 RedHawk 运行参数。

\begin{lstlisting}
                    -ct: specifies a file which conatins the LVS connection informations.
\end{lstlisting}

GSR 关键字 ：用于设置与  相关的 RedHawk 运行参数。

\begin{lstlisting}
                           CONNECT { LVSlayer1 LVSlayer2 }
                           CONNECT { LVSlayer1 LVSlayer2 BY LVSviaLayer }
                    -rcm: specifies the RC mapping file which contains the name mapping from LVS
\end{lstlisting}

层 to tech 文件 层 .

\begin{lstlisting}
                    -e: specifies the input EM iRCX file.
\end{lstlisting}

GSR 关键字 ：用于设置与  相关的 RedHawk 运行参数。

GSR 关键字 ：用于设置与  相关的 RedHawk 运行参数。（必需）

\begin{lstlisting}
            *.itf formats and translates the data into RedHawk readable tech file format. The syntax
\end{lstlisting}

GSR 关键字 ：用于设置与  相关的 RedHawk 运行参数。

\begin{lstlisting}
                  rhtech -i <input_file> -f <foundry_name> -n <process_node>
                  [-o <output_file>][-m <layer_mapping_file>] [-e <EM_file>]
                  [ -pe <EM_file>] [-t <temp_file>] -mim <MIMLayers> -enc
\end{lstlisting}

GSR 关键字 ：用于设置与  相关的 RedHawk 运行参数。

\begin{lstlisting}
                    -i <input_file> : specifies required vendor technology input file, in nxtgrd or itf
\end{lstlisting}

GSR 关键字 ：用于设置与  相关的 RedHawk 运行参数。

\begin{lstlisting}
                    -f <foundry_name> : includes the foundry name as a comment in generated Tech
\end{lstlisting}

GSR 关键字 ：用于设置与  相关的 RedHawk 运行参数。

\begin{lstlisting}
                    -n <process_node> : includes the process node information as a comment in
\end{lstlisting}

GSR 关键字 ：用于设置与  相关的 RedHawk 运行参数。

\begin{lstlisting}
                    -o <output_file> : specifies optional output file; default is '<input_file-
                        tech_name>.tech'. Note that the output file contains several comment lines
\end{lstlisting}

GSR 关键字 ：用于设置与  相关的 RedHawk 运行参数。

\begin{lstlisting}
                 # Apache Redhawk Technology File
                 # generated by 'rhtech' program from file- 'in/tddmy09m7t.itf'
                    -m <layer_mapping_file> : specifies optional layer name mapping file. No default.
                    -e <EM_file> : specifies basic EM current density limit file.
                    -pe <EM_file>: specifies special polynomial-based EM rule file. This will
\end{lstlisting}

GSR 关键字 ：用于设置与  相关的 RedHawk 运行参数。

\begin{lstlisting}
                 CONDUCTOR metal1 {
                      POLYNOMIAL_BASED_EM_DC {
                           LENGTH_RANGES { 3.402823e+38 }
                           EM_POLYNOMIAL { 4.14 * w }
                      }
                      POLYNOMIAL_BASED_EM_RMS {
                           LENGTH_RANGES { 3.402823e+38 }
                           EM_POLYNOMIAL { 3.64 * w *( sqrt ( 1 + ( 2.46 /( w + .08 ))))}
\end{lstlisting}

\section{gds2rh/gds2def}
\subsection*{gds2rh/gds2def}

从 GDSII 数据生成 DEF/LEF，用于 RedHawk 数据库集成。

GSR 关键字 ：用于设置与  相关的 RedHawk 运行参数。

\begin{lstlisting}
                          <via_layer_name> : name of via layer
                          <via_max_current>: maximum allowable current for a via (mA) whose area
\end{lstlisting}

GSR 关键字 ：用于设置与  相关的 RedHawk 运行参数。

\begin{lstlisting}
                           ---------------------------------------------
\end{lstlisting}

GSR 关键字 ：用于设置与  相关的 RedHawk 运行参数。

GSR 关键字 ：用于设置与  相关的 RedHawk 运行参数。

\begin{lstlisting}
                           -------------------------
\end{lstlisting}

过孔1 0.189 per 过孔 过孔2 0.210 per 过孔

GSR 关键字 ：用于设置与  相关的 RedHawk 运行参数。

\begin{lstlisting}
                    -t : Specifies the file providing the temperature at which the nxtgrd parameters
\end{lstlisting}

GSR 关键字 ：用于设置与  相关的 RedHawk 运行参数。

\begin{lstlisting}
                            <layer_name> <nominal_T> <RH_actual_T> <Tc1> <Tc2>
\end{lstlisting}

GSR 关键字 ：用于设置与  相关的 RedHawk 运行参数。

\begin{lstlisting}
                          <nominal_T> <RH_actual_T> : nominal and RedHawk actual temperature
\end{lstlisting}

GSR 关键字 ：用于设置与  相关的 RedHawk 运行参数。

\begin{lstlisting}
                          <Tc1> <Tc2> : first and second order temperature coefficients.
\end{lstlisting}

GSR 关键字 ：用于设置与  相关的 RedHawk 运行参数。

\begin{lstlisting}
                           <layer_name> 25 125 <Tc1> <Tc2>
\end{lstlisting}

GSR 关键字 ：用于设置与  相关的 RedHawk 运行参数。

\begin{lstlisting}
                    -enc: automatically inserts partial encryption flags #ENCRYPT_START/
                        #ENCRYPT_END to the non-EM part.
\end{lstlisting}

GSR 关键字 ：用于设置与  相关的 RedHawk 运行参数。

GSR 关键字 ：用于设置与  相关的 RedHawk 运行参数。

GSR 关键字 ：用于设置与  相关的 RedHawk 运行参数。

\subsection{GDS2RH/GDS2DEF 配置文件关键字}
\subsubsection*{GDS2RH/GDS2DEF Configuration File Keywords}

以下逐条译出 gds2rh/gds2def 配置文件中的主要关键字（语法、默认值与示例）。

\subsection{GDSII Files 关键字}
\subsubsection*{GDSII Files Keywords}

\subsubsection{\texttt{GDS_FILE}}
\paragraph*{GDS_FILE}

GSR 关键字 GDS_FILE：用于设置与 GDS、文件 相关的 RedHawk 运行参数。（必需）

\subsubsection{\texttt{GDS_MAP_FILE}}
\paragraph*{GDS_MAP_FILE}

GSR 关键字 GDS_MAP_FILE：用于设置与 GDS、MAP、文件 相关的 RedHawk 运行参数。

\textbf{语法：}
\begin{lstlisting}
GDS_MAP_FILE <GDSII_layer_map_file>
\end{lstlisting}

\subsection{Top Cell Definition 关键字}
\subsubsection*{Top Cell Definition Keywords}

\subsubsection{\texttt{CELL_NAME_MAP}}
\paragraph*{CELL_NAME_MAP}

GSR 关键字 CELL_NAME_MAP：用于设置与 单元、NAME、MAP 相关的 RedHawk 运行参数。

\textbf{语法：}
\begin{lstlisting}
CELL_NAME_MAP {
<gds_name> <desired name>
}
\end{lstlisting}

\subsubsection{\texttt{FLATTEN_TOP}}
\paragraph*{FLATTEN_TOP}

GSR 关键字 FLATTEN_TOP：用于设置与 FLATTEN、顶层 相关的 RedHawk 运行参数。（默认：0)）

\textbf{语法：}
\begin{lstlisting}
FLATTEN_TOP [0|1]
gds2rh/gds2def
\end{lstlisting}

\subsubsection{\texttt{FLATTEN_TOP_PROPERTIES}}
\paragraph*{FLATTEN_TOP_PROPERTIES}

GSR 关键字 FLATTEN_TOP_PROPERTIES：用于设置与 FLATTEN、顶层、PROPERTIES 相关的 RedHawk 运行参数。（默认：is 0, in which no optimized tracing is p）

\textbf{语法：}
\begin{lstlisting}
FLATTEN_TOP_PROPERTIES {
OPTIMIZED_TRACING [ 0 |1 |2 ]
}
\end{lstlisting}

\subsubsection{\texttt{GENERATE_TOP_LEF}}
\paragraph*{GENERATE_TOP_LEF}

GSR 关键字 GENERATE_TOP_LEF：用于设置与 GENERATE、顶层、LEF 相关的 RedHawk 运行参数。（默认：0）

\textbf{语法：}
\begin{lstlisting}
GENERATE_TOP_LEF [0 |1]
\end{lstlisting}

\subsubsection{\texttt{TOP_CELL}}
\paragraph*{TOP_CELL}

GSR 关键字 TOP_CELL：用于设置与 顶层、单元 相关的 RedHawk 运行参数。

\subsubsection{\texttt{TOP_CELLS}}
\paragraph*{TOP_CELLS}

GSR 关键字 TOP_CELLS：用于设置与 顶层、单元 相关的 RedHawk 运行参数。

\textbf{参数说明：}
\begin{itemize}
\item \texttt{<TopCell\_name>} — GSR 关键字 TOP\_CELLS：用于设置与 顶层、单元 相关的 RedHawk 运行参数。（可选）
\end{itemize}

\subsubsection{\texttt{TOP_CELL_MAP}}
\paragraph*{TOP_CELL_MAP}

GSR 关键字 TOP_CELL_MAP：用于设置与 顶层、单元、MAP 相关的 RedHawk 运行参数。

\subsubsection{\texttt{TOP_CELLS_MAP}}
\paragraph*{TOP_CELLS_MAP}

GSR 关键字 TOP_CELLS_MAP：用于设置与 顶层、单元、MAP 相关的 RedHawk 运行参数。

\subsubsection{\texttt{DEF_FILE_DEFINITIONS}}
\paragraph*{DEF_FILE_DEFINITIONS}

GSR 关键字 DEF_FILE_DEFINITIONS：用于设置与 DEF、文件、DEFINITIONS 相关的 RedHawk 运行参数。（默认：resolution units (default 2000) and die）

\textbf{语法：}
\begin{lstlisting}
DEF_FILE_DEFINITIONS {
UNITS <DEF resolution unit>
DIEAREA <llx> <lly> <urx> <ury>
gds2rh/gds2def
OFFSET <x_off> <y_off>
}
\end{lstlisting}

\subsection{Nets Definition 关键字}
\subsubsection*{Nets Definition Keywords}

\subsubsection{\texttt{CHECK_TRACING}}
\paragraph*{CHECK_TRACING}

GSR 关键字 CHECK_TRACING：用于设置与 检查、TRACING 相关的 RedHawk 运行参数。（默认：is 0 (do not check for CHECK_TRACING [0）

\subsubsection{\texttt{DUMMY_NET_OUTPUT_PINS_FILE}}
\paragraph*{DUMMY_NET_OUTPUT_PINS_FILE}

GSR 关键字 DUMMY_NET_OUTPUT_PINS_FILE：用于设置与 DUMMY、网络、OUTPUT、引脚、文件 相关的 RedHawk 运行参数。（必需）

\textbf{语法：}
\begin{lstlisting}
DUMMY_NET_OUTPUT_PINS_FILE <filename>
Sample output of NETS in DEF :
NETS 1 ;
- dummy_net (*Q)(*Qout) ...
\end{lstlisting}

\subsubsection{\texttt{EXTRACT_ALL_LEF_PIN_SIGNALS}}
\paragraph*{EXTRACT_ALL_LEF_PIN_SIGNALS}

GSR 关键字 EXTRACT_ALL_LEF_PIN_SIGNALS：用于设置与 提取、ALL、LEF、引脚、SIGNALS 相关的 RedHawk 运行参数。（默认：behavior）

\subsubsection{\texttt{GND_NETS}}
\paragraph*{GND_NETS}

GSR 关键字 GND_NETS：用于设置与 GND、网络 相关的 RedHawk 运行参数。

\subsubsection{\texttt{VDD_NETS}}
\paragraph*{VDD_NETS}

GSR 关键字 VDD_NETS：用于设置与 VDD、网络 相关的 RedHawk 运行参数。

\textbf{参数说明：}
\begin{itemize}
\item \texttt{<power/ground\_net\_name\_in\_output\_DEF>} — GSR 关键字 VDD\_NETS：用于设置与 VDD、网络 相关的 RedHawk 运行参数。
\end{itemize}

\textbf{示例：}
\begin{lstlisting}
example, if a power net only has one entry and the power net name to be used in output
DEF is the same as the net name in the entry, then a simplified format can be used, as
shown below:
VDD_NETS {
<power net1>
<power net2>
...
}
An example power net section is shown below:
VDD_NETS {
gds2rh/gds2def
VDD {
\end{lstlisting}

\subsubsection{\texttt{VDD}}
\paragraph*{VDD}

GSR 关键字 VDD：用于设置与 VDD 相关的 RedHawk 运行参数。

\textbf{示例：}
\begin{lstlisting}
example, VDD, VDD1_core, VDD2_core, and VDD3_core is included under VDD in the
generated .def file. Note that the power net names are case sensitive.
\end{lstlisting}

\subsubsection{\texttt{IGNORE_LEF_PG_PIN_GEOMETRY}}
\paragraph*{IGNORE_LEF_PG_PIN_GEOMETRY}

在数据输入或分析流程中忽略与 LEF、PG、引脚、GEOMETRY 相关的对象或检查。

\subsubsection{\texttt{MULTI_TASKS}}
\paragraph*{MULTI_TASKS}

GSR 关键字 MULTI_TASKS：用于设置与 MULTI、TASKS 相关的 RedHawk 运行参数。（默认：0）

\subsubsection{\texttt{NET_NAME_CASE_SENSITIVE}}
\paragraph*{NET_NAME_CASE_SENSITIVE}

GSR 关键字 NET_NAME_CASE_SENSITIVE：用于设置与 网络、NAME、CASE、SENSITIVE 相关的 RedHawk 运行参数。（默认：, text label name searching is case sens）

\subsubsection{\texttt{SIGNAL_NETS}}
\paragraph*{SIGNAL_NETS}

GSR 关键字 SIGNAL_NETS：用于设置与 SIGNAL、网络 相关的 RedHawk 运行参数。

\textbf{参数说明：}
\begin{itemize}
\item \texttt{<DEF\_signal\_net\_name-N>} — GSR 关键字 SIGNAL\_NETS：用于设置与 SIGNAL、网络 相关的 RedHawk 运行参数。
\end{itemize}

\subsubsection{\texttt{USE_TEXT_FROM_HIERARCHY_LEVEL}}
\paragraph*{USE_TEXT_FROM_HIERARCHY_LEVEL}

GSR 关键字 USE_TEXT_FROM_HIERARCHY_LEVEL：用于设置与 USE、TEXT、从、HIERARCHY、LEVEL 相关的 RedHawk 运行参数。

\textbf{语法：}
\begin{lstlisting}
USE_TEXT_FROM_HIERARCHY_LEVEL <hierarchy_level_descripton>
\end{lstlisting}

\textbf{示例：}
\begin{lstlisting}
USE_TEXT_FROM_HIERARCHY_LEVEL 0 2 5 ... # 0 is top-level
USE_TEXT_FROM_HIERARCHY_LEVEL 2-5
USE_TEXT_FROM_HIERARCHY_LEVEL 2-MAX
USE_TEXT_FROM_HIERARCHY_LEVEL 2 - MAX
USE_TEXT_FROM_HIERARCHY_LEVEL 0-5
\end{lstlisting}

\subsubsection{\texttt{USE_TEXT_FROM_CELLS}}
\paragraph*{USE_TEXT_FROM_CELLS}

GSR 关键字 USE_TEXT_FROM_CELLS：用于设置与 USE、TEXT、从、单元 相关的 RedHawk 运行参数。

\subsubsection{\texttt{USE_TEXT_FROM_CELLS}}
\paragraph*{USE_TEXT_FROM_CELLS}

GSR 关键字 USE_TEXT_FROM_CELLS：用于设置与 USE、TEXT、从、单元 相关的 RedHawk 运行参数。

\subsubsection{\texttt{USE_TEXT_FROM_CELLS}}
\paragraph*{USE_TEXT_FROM_CELLS}

GSR 关键字 USE_TEXT_FROM_CELLS：用于设置与 USE、TEXT、从、单元 相关的 RedHawk 运行参数。

\subsubsection{\texttt{USE_TEXT_FROM_HIERARCHY_BLOCK}}
\paragraph*{USE_TEXT_FROM_HIERARCHY_BLOCK}

GSR 关键字 USE_TEXT_FROM_HIERARCHY_BLOCK：用于设置与 USE、TEXT、从、HIERARCHY、模块 相关的 RedHawk 运行参数。

\textbf{语法：}
\begin{lstlisting}
USE_TEXT_FROM_HIERARCHY_BLOCK {
<hierarchy_path1>
...
}
\end{lstlisting}

\textbf{示例：}
\begin{lstlisting}
USE_TEXT_FROM_HIERARCHY_BLOCK {
TOP/A
TOP/A/B
\end{lstlisting}

\subsubsection{\texttt{TOP}}
\paragraph*{TOP}

GSR 关键字 TOP：用于设置与 顶层 相关的 RedHawk 运行参数。

\subsubsection{\texttt{USE_TOP_LEVEL_TEXT_ONLY}}
\paragraph*{USE_TOP_LEVEL_TEXT_ONLY}

GSR 关键字 USE_TOP_LEVEL_TEXT_ONLY：用于设置与 USE、顶层、LEVEL、TEXT、仅 相关的 RedHawk 运行参数。（默认：is 0）

\subsection{Layer Map Definition 关键字}
\subsubsection*{Layer Map Definition Keywords}

\subsubsection{\texttt{GDS_MAP_FILE}}
\paragraph*{GDS_MAP_FILE}

GSR 关键字 GDS_MAP_FILE：用于设置与 GDS、MAP、文件 相关的 RedHawk 运行参数。（必需）

\subsubsection{\texttt{DATATYPE}}
\paragraph*{DATATYPE}

GSR 关键字 DATATYPE：用于设置与 DATATYPE 相关的 RedHawk 运行参数。（可选）

\subsection{Input LEF 关键字}
\subsubsection*{Input LEF Keywords}

\subsubsection{\texttt{IGNORE_LEF_ORIGIN_FOR_START_PT}}
\paragraph*{IGNORE_LEF_ORIGIN_FOR_START_PT}

在数据输入或分析流程中忽略与 LEF、ORIGIN、FOR、START、PT 相关的对象或检查。（默认：is 0）

\subsubsection{\texttt{LEF_FILE}}
\paragraph*{LEF_FILE}

GSR 关键字 LEF_FILE：用于设置与 LEF、文件 相关的 RedHawk 运行参数。

\subsubsection{\texttt{BOX_CELLS_DEBUG}}
\paragraph*{BOX_CELLS_DEBUG}

GSR 关键字 BOX_CELLS_DEBUG：用于设置与 BOX、单元、DEBUG 相关的 RedHawk 运行参数。

\subsubsection{\texttt{LEF_FOREIGN_OFFSET_HANDLE}}
\paragraph*{LEF_FOREIGN_OFFSET_HANDLE}

GSR 关键字 LEF_FOREIGN_OFFSET_HANDLE：用于设置与 LEF、FOREIGN、OFFSET、HANDLE 相关的 RedHawk 运行参数。（默认：is 1）

\subsubsection{\texttt{USE_LEF_FOR_LOGICAL_CONNECTION}}
\paragraph*{USE_LEF_FOR_LOGICAL_CONNECTION}

GSR 关键字 USE_LEF_FOR_LOGICAL_CONNECTION：用于设置与 USE、LEF、FOR、LOGICAL、连接 相关的 RedHawk 运行参数。

\subsubsection{\texttt{USE_LEF_PINS_FOR_TRACING}}
\paragraph*{USE_LEF_PINS_FOR_TRACING}

GSR 关键字 USE_LEF_PINS_FOR_TRACING：用于设置与 USE、LEF、引脚、FOR、TRACING 相关的 RedHawk 运行参数。

\textbf{参数说明：}
GSR 关键字 USE\_LEF\_PINS\_FOR\_TRACING：用于设置与 USE、LEF、引脚、FOR、TRACING 相关的 RedHawk 运行参数。（默认：0 (off, no LEF pin based pin tracing)）

\subsubsection{\texttt{USE_LEF_FOR_TOP_BBOX}}
\paragraph*{USE_LEF_FOR_TOP_BBOX}

GSR 关键字 USE_LEF_FOR_TOP_BBOX：用于设置与 USE、LEF、FOR、顶层、BBOX 相关的 RedHawk 运行参数。（默认：(1), forces extraction to pick up the ce）

\textbf{语法：}
\begin{lstlisting}
USE_LEF_FOR_TOP_BBOX [0|1]
\end{lstlisting}

\subsubsection{\texttt{LEF_PIN_POWER}}
\paragraph*{LEF_PIN_POWER}

GSR 关键字 LEF_PIN_POWER：用于设置与 LEF、引脚、电源 相关的 RedHawk 运行参数。（默认：0 (off) LEF_PIN_POWER [0|1]）

\subsubsection{\texttt{GENERATE_SEPARATE_PG_PINS_LEF}}
\paragraph*{GENERATE_SEPARATE_PG_PINS_LEF}

GSR 关键字 GENERATE_SEPARATE_PG_PINS_LEF：用于设置与 GENERATE、SEPARATE、PG、引脚、LEF 相关的 RedHawk 运行参数。（默认：is off (0)）

\subsubsection{\texttt{IGNORE_LEF_FOREIGN_OFFSET}}
\paragraph*{IGNORE_LEF_FOREIGN_OFFSET}

在数据输入或分析流程中忽略与 LEF、FOREIGN、OFFSET 相关的对象或检查。（默认：is 0 or off, which means the FOREIGN off） 注意：在数据输入或分析流程中忽略与 LEF、FOREIGN、OFFSET 相关的对象或检查。

\begin{noteBox}
在数据输入或分析流程中忽略与 LEF、FOREIGN、OFFSET 相关的对象或检查。
\end{noteBox}

\subsection{Geometry Extraction 关键字}
\subsubsection*{Geometry Extraction Keywords}

\subsubsection{\texttt{APACHE_PHYSICAL_MODEL}}
\paragraph*{APACHE_PHYSICAL_MODEL}

GSR 关键字 APACHE_PHYSICAL_MODEL：用于设置与 APACHE、PHYSICAL、MODEL 相关的 RedHawk 运行参数。（默认：none）

\textbf{语法：}
\begin{lstlisting}
APACHE_PHYSICAL_MODEL <APM_metal_layer>
\end{lstlisting}

\subsubsection{\texttt{ADJUST_POLYGON}}
\paragraph*{ADJUST_POLYGON}

GSR 关键字 ADJUST_POLYGON：用于设置与 ADJUST、POLYGON 相关的 RedHawk 运行参数。（默认：is 0）

\textbf{语法：}
\begin{lstlisting}
ADJUST_POLYGON [0 |1]
\end{lstlisting}

\subsubsection{\texttt{AUTOFIT_DIFFUSION_WITH_POLY}}
\paragraph*{AUTOFIT_DIFFUSION_WITH_POLY}

GSR 关键字 AUTOFIT_DIFFUSION_WITH_POLY：用于设置与 AUTOFIT、DIFFUSION、WITH、POLY 相关的 RedHawk 运行参数。

\textbf{语法：}
\begin{lstlisting}
AUTOFIT_DIFFUSION_WITH_POLY [0 |1]
\end{lstlisting}

\subsubsection{\texttt{ENABLE_POLY_VIA_MERGE}}
\paragraph*{ENABLE_POLY_VIA_MERGE}

启用 POLY、过孔、MERGE 相关功能或检查。（默认：Off (0)）

\textbf{语法：}
\begin{lstlisting}
ENABLE_POLY_VIA_MERGE [ 0 | 1 ]
\end{lstlisting}

\subsubsection{\texttt{EXTRACTION_STARTING_LAYER}}
\paragraph*{EXTRACTION_STARTING_LAYER}

GSR 关键字 EXTRACTION_STARTING_LAYER：用于设置与 提取、STARTING、层 相关的 RedHawk 运行参数。（默认：is to trace and extract all the layers s）

\subsubsection{\texttt{EXTRACT_ISOLATED_SUBDOMAINS}}
\paragraph*{EXTRACT_ISOLATED_SUBDOMAINS}

GSR 关键字 EXTRACT_ISOLATED_SUBDOMAINS：用于设置与 提取、ISOLATED、SUBDOMAINS 相关的 RedHawk 运行参数。

\subsubsection{\texttt{IMPORT_REGIONS}}
\paragraph*{IMPORT_REGIONS}

GSR 关键字 IMPORT_REGIONS：用于设置与 导入、REGIONS 相关的 RedHawk 运行参数。

\subsubsection{\texttt{MERGE_WIRE_NETWISE}}
\paragraph*{MERGE_WIRE_NETWISE}

GSR 关键字 MERGE_WIRE_NETWISE：用于设置与 MERGE、WIRE、NETWISE 相关的 RedHawk 运行参数。（默认：for gds2rh -m design flows, and 0 by def）

\subsubsection{\texttt{NEW_IMPORT_REGION}}
\paragraph*{NEW_IMPORT_REGION}

GSR 关键字 NEW_IMPORT_REGION：用于设置与 NEW、导入、区域 相关的 RedHawk 运行参数。（默认：Off；必需）

\subsubsection{\texttt{NO_METAL_IN_VIA_MODEL}}
\paragraph*{NO_METAL_IN_VIA_MODEL}

GSR 关键字 NO_METAL_IN_VIA_MODEL：用于设置与 NO、金属、IN、过孔、MODEL 相关的 RedHawk 运行参数。（默认：(0), via models generated are the same s）

\subsubsection{\texttt{NON_RECT_DIEAREA}}
\paragraph*{NON_RECT_DIEAREA}

GSR 关键字 NON_RECT_DIEAREA：用于设置与 NON、RECT、DIEAREA 相关的 RedHawk 运行参数。

\subsubsection{\texttt{OUTPUT_VIA_ARRAY}}
\paragraph*{OUTPUT_VIA_ARRAY}

GSR 关键字 OUTPUT_VIA_ARRAY：用于设置与 OUTPUT、过孔、ARRAY 相关的 RedHawk 运行参数。

\subsubsection{\texttt{START_POINT_TOLERANCE}}
\paragraph*{START_POINT_TOLERANCE}

GSR 关键字 START_POINT_TOLERANCE：用于设置与 START、POINT、TOLERANCE 相关的 RedHawk 运行参数。（可选；默认：behavior of the algorithm attempts to au）

\textbf{语法：}
\begin{lstlisting}
START_POINT_TOLERANCE <search_radius>
\end{lstlisting}

\textbf{示例：}
\begin{lstlisting}
START_POINT_TOLERANCE 0.05
\end{lstlisting}

\subsubsection{\texttt{VIA_NAME_PREFIX}}
\paragraph*{VIA_NAME_PREFIX}

GSR 关键字 VIA_NAME_PREFIX：用于设置与 过孔、NAME、PREFIX 相关的 RedHawk 运行参数。（默认：), a via model name could be GDSVia_logi）

\subsection{Selective Cell Hierarchy Handling 关键字}
\subsubsection*{Selective Cell Hierarchy Handling Keywords}

\subsubsection{\texttt{WHITE_BOX_CELLS}}
\paragraph*{WHITE_BOX_CELLS}

GSR 关键字 WHITE_BOX_CELLS：用于设置与 WHITE、BOX、单元 相关的 RedHawk 运行参数。 注意：GSR 关键字 WHITE_BOX_CELLS：用于设置与 WHITE、BOX、单元 相关的 RedHawk 运行参数。

\begin{noteBox}
GSR 关键字 WHITE\_BOX\_CELLS：用于设置与 WHITE、BOX、单元 相关的 RedHawk 运行参数。
\end{noteBox}

\subsubsection{\texttt{BLACK_BOX_CELLS}}
\paragraph*{BLACK_BOX_CELLS}

GSR 关键字 BLACK_BOX_CELLS：用于设置与 BLACK、BOX、单元 相关的 RedHawk 运行参数。

\subsubsection{\texttt{GRAY_BOX_CELLS}}
\paragraph*{GRAY_BOX_CELLS}

GSR 关键字 GRAY_BOX_CELLS：用于设置与 GRAY、BOX、单元 相关的 RedHawk 运行参数。（默认：0）

\subsubsection{\texttt{BOX_CELLS_PROPERTIES}}
\paragraph*{BOX_CELLS_PROPERTIES}

GSR 关键字 BOX_CELLS_PROPERTIES：用于设置与 BOX、单元、PROPERTIES 相关的 RedHawk 运行参数。（默认：behavior when BOX_CELLS_PROPERTIES is no）

\subsubsection{\texttt{CREATE_LEF_MACRO_FOR_BOX_CELLS}}
\paragraph*{CREATE_LEF_MACRO_FOR_BOX_CELLS}

GSR 关键字 CREATE_LEF_MACRO_FOR_BOX_CELLS：用于设置与 CREATE、LEF、宏、FOR、BOX、单元 相关的 RedHawk 运行参数。（默认：is 0 (no LEF file)）

\textbf{参数说明：}
\begin{itemize}
\item \texttt{1} — GSR 关键字 CREATE\_LEF\_MACRO\_FOR\_BOX\_CELLS：用于设置与 CREATE、LEF、宏、FOR、BOX、单元 相关的 RedHawk 运行参数。
\end{itemize}

\subsubsection{\texttt{UNIQUIFY_BOX_CELLS}}
\paragraph*{UNIQUIFY_BOX_CELLS}

GSR 关键字 UNIQUIFY_BOX_CELLS：用于设置与 UNIQUIFY、BOX、单元 相关的 RedHawk 运行参数。（默认：'2' value assigns a LEF pin name based o）

\subsection{Auto Pad Location Generation 关键字}
\subsubsection*{Auto Pad Location Generation Keywords}

\subsubsection{\texttt{GENERATE_PLOC}}
\paragraph*{GENERATE_PLOC}

GSR 关键字 GENERATE_PLOC：用于设置与 GENERATE、焊盘位置 相关的 RedHawk 运行参数。

\textbf{参数说明：}
\begin{itemize}
\item \texttt{USE\_TEXT\_LABEL} — GSR 关键字 GENERATE\_PLOC：用于设置与 GENERATE、焊盘位置 相关的 RedHawk 运行参数。
\end{itemize}

\subsubsection{\texttt{SIG}}
\paragraph*{SIG}

GSR 关键字 SIG：用于设置与 SIG 相关的 RedHawk 运行参数。 注意：GSR 关键字 SIG：用于设置与 SIG 相关的 RedHawk 运行参数。

\begin{noteBox}
GSR 关键字 SIG：用于设置与 SIG 相关的 RedHawk 运行参数。
\end{noteBox}

\subsubsection{\texttt{GENERATE_PLOC_FILTER}}
\paragraph*{GENERATE_PLOC_FILTER}

GSR 关键字 GENERATE_PLOC_FILTER：用于设置与 GENERATE、焊盘位置、FILTER 相关的 RedHawk 运行参数。

\textbf{参数说明：}
GSR 关键字 GENERATE\_PLOC\_FILTER：用于设置与 GENERATE、焊盘位置、FILTER 相关的 RedHawk 运行参数。

\subsection{DSPF/SPEF-based standard cell flow 关键字}
\subsubsection*{DSPF/SPEF-based standard cell flow Keywords}

\subsubsection{\texttt{SPEF_FILE}}
\paragraph*{SPEF_FILE}

GSR 关键字 SPEF_FILE：用于设置与 SPEF、文件 相关的 RedHawk 运行参数。

\subsubsection{\texttt{SPEF_CELLS}}
\paragraph*{SPEF_CELLS}

GSR 关键字 SPEF_CELLS：用于设置与 SPEF、单元 相关的 RedHawk 运行参数。

\subsubsection{\texttt{DEF_FILE}}
\paragraph*{DEF_FILE}

GSR 关键字 DEF_FILE：用于设置与 DEF、文件 相关的 RedHawk 运行参数。

\subsubsection{\texttt{GRAY_BOX_CELLS}}
\paragraph*{GRAY_BOX_CELLS}

GSR 关键字 GRAY_BOX_CELLS：用于设置与 GRAY、BOX、单元 相关的 RedHawk 运行参数。

\subsubsection{\texttt{DSPF_CELLS_PREFIX}}
\paragraph*{DSPF_CELLS_PREFIX}

GSR 关键字 DSPF_CELLS_PREFIX：用于设置与 DSPF、单元、PREFIX 相关的 RedHawk 运行参数。

\subsubsection{\texttt{SPEF_CELLS_PREFIX}}
\paragraph*{SPEF_CELLS_PREFIX}

GSR 关键字 SPEF_CELLS_PREFIX：用于设置与 SPEF、单元、PREFIX 相关的 RedHawk 运行参数。

\textbf{语法：}
\begin{lstlisting}
DSPF_CELLS_PREFIX {
<Block_prefix1>
...
}
For example, suppose ‘Top’ has three blocks, BlockA, BlockB and BlockC, and standard
cell buffer “BUFX1” is in all three blocks. Then the LEF includes “MACRO BUFX1“, and
GDS includes the following:
TOP    : Master Cell BUFX1
BlockA : Master Cell BlockA_BUFX1
BlockB : Master Cell BlockB_BUFX1
BlockC : Master Cell BlockC_BUFX1
\end{lstlisting}

\subsubsection{\texttt{NO_DEF_LOGICAL_CONNECTION}}
\paragraph*{NO_DEF_LOGICAL_CONNECTION}

GSR 关键字 NO_DEF_LOGICAL_CONNECTION：用于设置与 NO、DEF、LOGICAL、连接 相关的 RedHawk 运行参数。

\subsubsection{\texttt{DSPF_CELL_INSTANCE_MATCH_TOLERANCE}}
\paragraph*{DSPF_CELL_INSTANCE_MATCH_TOLERANCE}

GSR 关键字 DSPF_CELL_INSTANCE_MATCH_TOLERANCE：用于设置与 DSPF、单元、实例、MATCH、TOLERANCE 相关的 RedHawk 运行参数。（默认：tolerance is 0）

\subsubsection{\texttt{SPEF_MAGNIFICATION_FACTOR}}
\paragraph*{SPEF_MAGNIFICATION_FACTOR}

GSR 关键字 SPEF_MAGNIFICATION_FACTOR：用于设置与 SPEF、MAGNIFICATION、因子 相关的 RedHawk 运行参数。（默认：is 1）

\subsubsection{\texttt{USE_LEF_FOR_LOGICAL_CONNECTION}}
\paragraph*{USE_LEF_FOR_LOGICAL_CONNECTION}

GSR 关键字 USE_LEF_FOR_LOGICAL_CONNECTION：用于设置与 USE、LEF、FOR、LOGICAL、连接 相关的 RedHawk 运行参数。

\subsection{Automated Switch Cell Handling 关键字}
\subsubsection*{Automated Switch Cell Handling Keywords}

\subsubsection{\texttt{DEFINE_SWITCH_CELLS}}
\paragraph*{DEFINE_SWITCH_CELLS}

GSR 关键字 DEFINE_SWITCH_CELLS：用于设置与 定义、SWITCH、单元 相关的 RedHawk 运行参数。

\textbf{示例：}
\begin{lstlisting}
Example
DEFINE_SWITCH_CELLS {
sw1 HEADER VDD_EXT VDD_INT EN
sw2 HEADER VDD_EXT VDD_INT EN
}
\end{lstlisting}

\subsubsection{\texttt{EXTRACT_SWITCH_CELLS}}
\paragraph*{EXTRACT_SWITCH_CELLS}

GSR 关键字 EXTRACT_SWITCH_CELLS：用于设置与 提取、SWITCH、单元 相关的 RedHawk 运行参数。

\textbf{示例：}
\begin{lstlisting}
Example
EXTRACT_SWITCH_CELLS {
sw1 HEADER ext_vdd1 int_vdd1
sw1 HEADER ext_vdd2 int_vdd2
}
\end{lstlisting}

\subsubsection{\texttt{SWITCH_CELLS_PROPERTIES}}
\paragraph*{SWITCH_CELLS_PROPERTIES}

GSR 关键字 SWITCH_CELLS_PROPERTIES：用于设置与 SWITCH、单元、PROPERTIES 相关的 RedHawk 运行参数。

\subsubsection{\texttt{USE_SWITCH_XTOR}}
\paragraph*{USE_SWITCH_XTOR}

GSR 关键字 USE_SWITCH_XTOR：用于设置与 USE、SWITCH、XTOR 相关的 RedHawk 运行参数。

\subsubsection{\texttt{XTOR_FINGER_PITCH}}
\paragraph*{XTOR_FINGER_PITCH}

GSR 关键字 XTOR_FINGER_PITCH：用于设置与 XTOR、FINGER、PITCH 相关的 RedHawk 运行参数。

\subsubsection{\texttt{XTOR_FINGER_WIDTH}}
\paragraph*{XTOR_FINGER_WIDTH}

GSR 关键字 XTOR_FINGER_WIDTH：用于设置与 XTOR、FINGER、WIDTH 相关的 RedHawk 运行参数。

\subsubsection{\texttt{EXTEND_SWITCH_BBOX}}
\paragraph*{EXTEND_SWITCH_BBOX}

GSR 关键字 EXTEND_SWITCH_BBOX：用于设置与 EXTEND、SWITCH、BBOX 相关的 RedHawk 运行参数。

\subsubsection{\texttt{MAX_SWITCH_WIDTH}}
\paragraph*{MAX_SWITCH_WIDTH}

GSR 关键字 MAX_SWITCH_WIDTH：用于设置与 MAX、SWITCH、WIDTH 相关的 RedHawk 运行参数。

\subsubsection{\texttt{MAX_SWITCH_HEIGHT}}
\paragraph*{MAX_SWITCH_HEIGHT}

GSR 关键字 MAX_SWITCH_HEIGHT：用于设置与 MAX、SWITCH、HEIGHT 相关的 RedHawk 运行参数。

\textbf{参数说明：}
GSR 关键字 MAX\_SWITCH\_HEIGHT：用于设置与 MAX、SWITCH、HEIGHT 相关的 RedHawk 运行参数。（默认：=0 MAX\_SWITCH\_WIDTH - to restrict the ma）

\subsubsection{\texttt{USE_SWITCH_XTOR}}
\paragraph*{USE_SWITCH_XTOR}

GSR 关键字 USE_SWITCH_XTOR：用于设置与 USE、SWITCH、XTOR 相关的 RedHawk 运行参数。

\subsubsection{\texttt{SWITCH_MODEL_TABLE}}
\paragraph*{SWITCH_MODEL_TABLE}

GSR 关键字 SWITCH_MODEL_TABLE：用于设置与 SWITCH、MODEL、TABLE 相关的 RedHawk 运行参数。 注意：GSR 关键字 SWITCH_MODEL_TABLE：用于设置与 SWITCH、MODEL、TABLE 相关的 RedHawk 运行参数。

\begin{noteBox}
GSR 关键字 SWITCH\_MODEL\_TABLE：用于设置与 SWITCH、MODEL、TABLE 相关的 RedHawk 运行参数。
\end{noteBox}

\subsection{Other 关键字}
\subsubsection*{Other Keywords}

\subsubsection{\texttt{GENERATE_GSR_ONLY}}
\paragraph*{GENERATE_GSR_ONLY}

GSR 关键字 GENERATE_GSR_ONLY：用于设置与 GENERATE、GSR、仅 相关的 RedHawk 运行参数。

\subsubsection{\texttt{OUTPUT_DIRECTORY}}
\paragraph*{OUTPUT_DIRECTORY}

GSR 关键字 OUTPUT_DIRECTORY：用于设置与 OUTPUT、DIRECTORY 相关的 RedHawk 运行参数。

\subsubsection{\texttt{COMPRESS_DEF}}
\paragraph*{COMPRESS_DEF}

GSR 关键字 COMPRESS_DEF：用于设置与 COMPRESS、DEF 相关的 RedHawk 运行参数。（默认：value is '0' (not gzipped)）

\subsubsection{\texttt{BUSBITCHARS}}
\paragraph*{BUSBITCHARS}

GSR 关键字 BUSBITCHARS：用于设置与 BUSBITCHARS 相关的 RedHawk 运行参数。（默认：value “[ ]”）

\subsubsection{\texttt{DEFAULT_NET_CURRENT_PIN_CREATION}}
\paragraph*{DEFAULT_NET_CURRENT_PIN_CREATION}

GSR 关键字 DEFAULT_NET_CURRENT_PIN_CREATION：用于设置与 DEFAULT、网络、CURRENT、引脚、CREATION 相关的 RedHawk 运行参数。（默认：current pins are created on the EXTRACTI）

\subsubsection{\texttt{ESD_CELLS}}
\paragraph*{ESD_CELLS}

GSR 关键字 ESD_CELLS：用于设置与 ESD、单元 相关的 RedHawk 运行参数。 注意：GSR 关键字 ESD_CELLS：用于设置与 ESD、单元 相关的 RedHawk 运行参数。

\textbf{参数说明：}
\begin{itemize}
\item \texttt{USE\_CELL\_NAMES/ CELL\_NAMES} — GSR 关键字 ESD\_CELLS：用于设置与 ESD、单元 相关的 RedHawk 运行参数。（默认：value of 1 means that the Clamp Cell Hie）
\end{itemize}

\begin{noteBox}
GSR 关键字 ESD\_CELLS：用于设置与 ESD、单元 相关的 RedHawk 运行参数。
\end{noteBox}

\subsubsection{\texttt{EXTRACT_METAL_FILLER}}
\paragraph*{EXTRACT_METAL_FILLER}

GSR 关键字 EXTRACT_METAL_FILLER：用于设置与 提取、金属、FILLER 相关的 RedHawk 运行参数。（默认：0）

\subsubsection{\texttt{GDS_INSTANCE_NAME_KEY}}
\paragraph*{GDS_INSTANCE_NAME_KEY}

GSR 关键字 GDS_INSTANCE_NAME_KEY：用于设置与 GDS、实例、NAME、KEY 相关的 RedHawk 运行参数。

\subsubsection{\texttt{INTERNAL_DBUNIT}}
\paragraph*{INTERNAL_DBUNIT}

GSR 关键字 INTERNAL_DBUNIT：用于设置与 INTERNAL、DBUNIT 相关的 RedHawk 运行参数。（默认：is 2000）

\subsubsection{\texttt{MAGIC_COMPATIBLE}}
\paragraph*{MAGIC_COMPATIBLE}

GSR 关键字 MAGIC_COMPATIBLE：用于设置与 MAGIC、COMPATIBLE 相关的 RedHawk 运行参数。（默认：is 0）

\textbf{语法：}
\begin{lstlisting}
MAGIC_COMPATIBLE [0|1]
\end{lstlisting}

\subsubsection{\texttt{MERGE_DUP_STRUCTURE}}
\paragraph*{MERGE_DUP_STRUCTURE}

GSR 关键字 MERGE_DUP_STRUCTURE：用于设置与 MERGE、DUP、STRUCTURE 相关的 RedHawk 运行参数。（默认：, the most recently processed structure）

\subsubsection{\texttt{MEMORY_SAVING_MODE}}
\paragraph*{MEMORY_SAVING_MODE}

GSR 关键字 MEMORY_SAVING_MODE：用于设置与 MEMORY、SAVING、模式 相关的 RedHawk 运行参数。（默认：MEMORY_SAVING_MODE [ 0 |1 ]）

\subsubsection{\texttt{OUTPUT_APACHECELL}}
\paragraph*{OUTPUT_APACHECELL}

GSR 关键字 OUTPUT_APACHECELL：用于设置与 OUTPUT、APACHECELL 相关的 RedHawk 运行参数。（默认：is 1）

\subsubsection{\texttt{OUTPUT_DETAILED_DEF_PIN}}
\paragraph*{OUTPUT_DETAILED_DEF_PIN}

GSR 关键字 OUTPUT_DETAILED_DEF_PIN：用于设置与 OUTPUT、DETAILED、DEF、引脚 相关的 RedHawk 运行参数。（默认：1）

\subsubsection{\texttt{FULL_GDSII_INFO}}
\paragraph*{FULL_GDSII_INFO}

GSR 关键字 FULL_GDSII_INFO：用于设置与 FULL、GDSII、INFO 相关的 RedHawk 运行参数。（默认：, GDS2RH generates the imported GDS data）

\subsubsection{\texttt{PREFLATTEN_CELLS}}
\paragraph*{PREFLATTEN_CELLS}

GSR 关键字 PREFLATTEN_CELLS：用于设置与 PREFLATTEN、单元 相关的 RedHawk 运行参数。 注意：GSR 关键字 PREFLATTEN_CELLS：用于设置与 PREFLATTEN、单元 相关的 RedHawk 运行参数。

\begin{noteBox}
GSR 关键字 PREFLATTEN\_CELLS：用于设置与 PREFLATTEN、单元 相关的 RedHawk 运行参数。
\end{noteBox}

\subsubsection{\texttt{REGION_BASED_PRATIO}}
\paragraph*{REGION_BASED_PRATIO}

GSR 关键字 REGION_BASED_PRATIO：用于设置与 区域、BASED、PRATIO 相关的 RedHawk 运行参数。

\textbf{语法：}
\begin{lstlisting}
REGION_BASED_PRATIO {
<REGION NAME> ALL <xll> <yll> <xur> <yur> <pratio_value>
............
BIT_CELL_REGION ALL <AUTO / BITCELL_NAME> <pratio val>
gds2rh/gds2def
REST_REGION ALL <pratio_value>
}
\end{lstlisting}

\subsubsection{\texttt{VALIDATE_MODEL}}
\paragraph*{VALIDATE_MODEL}

GSR 关键字 VALIDATE_MODEL：用于设置与 VALIDATE、MODEL 相关的 RedHawk 运行参数。（可选；默认：none）

\textbf{语法：}
\begin{lstlisting}
VALIDATE_MODEL {
TECH_FILE <tech_file_path>
LIB_FILES {
<lib_files>
...
}
PAD_FILES {
<pad_files>
...
}
ADD_PLOC_FROM_TOP_DEF 1
}
Running gds2rh or gds2def
To run gds2rh or gds2def, use the following invocation syntax:
[gds2drh | gds2def ] <config_file> ? <GDSII_filename> ?
? -out_dir <output_dir> ? -log <log_dir> ?
RETAIN_PRBOUNDARY_FOR_TOP_BBOX [ 1 | 0]
When you use ‘RETAIN_PRBOUNDARY_FOR_BOX_CELL_BBOX 1’, GDS2DB honors
the sub-block PRBOUNDARY keyword to get the sub-block bounding box in the
hierarchical flow. This enables you to provide LEF for sub-blocks (generated again with
PRBOUNDARY in layer map), which can be stitched up with the top-level seamlessly.
RETAIN_PRBOUNDARY_FOR_BOX_CELL_BBOX [ 1 | 0]
Specific metal resistor support using the marker layer
By default GDS traces all metal layers. If you want to break GDS tracing, a metal resistor
on the marker layer can be used. For specific metal-resistor situations in which two nets,
\end{lstlisting}

\textbf{参数说明：}
\begin{itemize}
\item \texttt{<config\_file>} — GSR 关键字 VALIDATE\_MODEL：用于设置与 VALIDATE、MODEL 相关的 RedHawk 运行参数。（默认：log files are written to）
\end{itemize}

\subsubsection{\texttt{THROUGHVIA}}
\paragraph*{THROUGHVIA}

GSR 关键字 THROUGHVIA：用于设置与 THROUGHVIA 相关的 RedHawk 运行参数。

\textbf{参数说明：}
\begin{itemize}
\item \texttt{TSV} — GSR 关键字 THROUGHVIA：用于设置与 THROUGHVIA 相关的 RedHawk 运行参数。（必需）
\end{itemize}

\textbf{示例：}
\begin{lstlisting}
TSV tv 10 -
gds2rh/gds2def
BM m 11 -
VIA_DUM v 999 -
M1 m 31 131
MCAP c 2 - MCAPM2VIA
VIA1 v 51 -
M2 m 32 132
THROUGHVIA TSV BM M2
\end{lstlisting}

\subsubsection{\texttt{USE_LEF_FOR_LOGICAL_CONNECTION}}
\paragraph*{USE_LEF_FOR_LOGICAL_CONNECTION}

GSR 关键字 USE_LEF_FOR_LOGICAL_CONNECTION：用于设置与 USE、LEF、FOR、LOGICAL、连接 相关的 RedHawk 运行参数。

\subsubsection{\texttt{PPL_ANA_LOC_65}}
\paragraph*{PPL_ANA_LOC_65}

GSR 关键字 PPL_ANA_LOC_65：用于设置与 PPL、ANA、LOC、65 相关的 RedHawk 运行参数。

\subsubsection{\texttt{IVHVY4}}
\paragraph*{IVHVY4}

GSR 关键字 IVHVY4：用于设置与 IVHVY4 相关的 RedHawk 运行参数。

\subsubsection{\texttt{IP1_ANA_LOC_65}}
\paragraph*{IP1_ANA_LOC_65}

GSR 关键字 IP1_ANA_LOC_65：用于设置与 IP1、ANA、LOC、65 相关的 RedHawk 运行参数。

\subsubsection{\texttt{VDDIOCO_5V_ANA_LOC_58_1}}
\paragraph*{VDDIOCO_5V_ANA_LOC_58_1}

GSR 关键字 VDDIOCO_5V_ANA_LOC_58_1：用于设置与 VDDIOCO、5V、ANA、LOC、58、1 相关的 RedHawk 运行参数。

\subsubsection{\texttt{IP1_ANA_LOC_65}}
\paragraph*{IP1_ANA_LOC_65}

GSR 关键字 IP1_ANA_LOC_65：用于设置与 IP1、ANA、LOC、65 相关的 RedHawk 运行参数。

\subsubsection{\texttt{IVHVY4}}
\paragraph*{IVHVY4}

GSR 关键字 IVHVY4：用于设置与 IVHVY4 相关的 RedHawk 运行参数。

\subsection{Configuration 关键字}
\subsubsection*{Configuration Keywords}

\subsubsection{\texttt{VDD}}
\paragraph*{VDD}

GSR 关键字 VDD：用于设置与 VDD 相关的 RedHawk 运行参数。

\subsubsection{\texttt{VSS}}
\paragraph*{VSS}

GSR 关键字 VSS：用于设置与 VSS 相关的 RedHawk 运行参数。

\subsection{Switch Cell Handling 关键字}
\subsubsection*{Switch Cell Handling Keywords}

\subsubsection{\texttt{DEFINE_SWITCH_CELLS}}
\paragraph*{DEFINE_SWITCH_CELLS}

GSR 关键字 DEFINE_SWITCH_CELLS：用于设置与 定义、SWITCH、单元 相关的 RedHawk 运行参数。

\textbf{示例：}
\begin{lstlisting}
Example
DEFINE_SWITCH_CELLS {
sw1 HEADER VDD_EXT VDD_INT EN
sw2 HEADER VDD_EXT VDD_INT EN
}
\end{lstlisting}

\subsubsection{\texttt{EXTRACT_SWITCH_CELLS}}
\paragraph*{EXTRACT_SWITCH_CELLS}

GSR 关键字 EXTRACT_SWITCH_CELLS：用于设置与 提取、SWITCH、单元 相关的 RedHawk 运行参数。

\textbf{示例：}
\begin{lstlisting}
Example
EXTRACT_SWITCH_CELLS {
sw1 HEADER ext_vdd1 int_vdd1
sw1 HEADER ext_vdd2 int_vdd2
}
\end{lstlisting}

\subsubsection{\texttt{SWITCH_CELLS_PROPERTIES}}
\paragraph*{SWITCH_CELLS_PROPERTIES}

GSR 关键字 SWITCH_CELLS_PROPERTIES：用于设置与 SWITCH、单元、PROPERTIES 相关的 RedHawk 运行参数。

\subsubsection{\texttt{USE_SWITCH_XTOR}}
\paragraph*{USE_SWITCH_XTOR}

GSR 关键字 USE_SWITCH_XTOR：用于设置与 USE、SWITCH、XTOR 相关的 RedHawk 运行参数。

\subsubsection{\texttt{XTOR_FINGER_PITCH}}
\paragraph*{XTOR_FINGER_PITCH}

GSR 关键字 XTOR_FINGER_PITCH：用于设置与 XTOR、FINGER、PITCH 相关的 RedHawk 运行参数。

\subsubsection{\texttt{XTOR_FINGER_WIDTH}}
\paragraph*{XTOR_FINGER_WIDTH}

GSR 关键字 XTOR_FINGER_WIDTH：用于设置与 XTOR、FINGER、WIDTH 相关的 RedHawk 运行参数。

\subsubsection{\texttt{EXTEND_SWITCH_BBOX}}
\paragraph*{EXTEND_SWITCH_BBOX}

GSR 关键字 EXTEND_SWITCH_BBOX：用于设置与 EXTEND、SWITCH、BBOX 相关的 RedHawk 运行参数。

\subsubsection{\texttt{MAX_SWITCH_WIDTH}}
\paragraph*{MAX_SWITCH_WIDTH}

GSR 关键字 MAX_SWITCH_WIDTH：用于设置与 MAX、SWITCH、WIDTH 相关的 RedHawk 运行参数。

\subsubsection{\texttt{MAX_SWITCH_HEIGHT}}
\paragraph*{MAX_SWITCH_HEIGHT}

GSR 关键字 MAX_SWITCH_HEIGHT：用于设置与 MAX、SWITCH、HEIGHT 相关的 RedHawk 运行参数。

\textbf{参数说明：}
GSR 关键字 MAX\_SWITCH\_HEIGHT：用于设置与 MAX、SWITCH、HEIGHT 相关的 RedHawk 运行参数。（默认：=0 MAX\_SWITCH\_WIDTH - to restrict the ma）

\subsubsection{\texttt{USE_SWITCH_XTOR}}
\paragraph*{USE_SWITCH_XTOR}

GSR 关键字 USE_SWITCH_XTOR：用于设置与 USE、SWITCH、XTOR 相关的 RedHawk 运行参数。

\subsubsection{\texttt{SWITCH_MODEL_TABLE}}
\paragraph*{SWITCH_MODEL_TABLE}

GSR 关键字 SWITCH_MODEL_TABLE：用于设置与 SWITCH、MODEL、TABLE 相关的 RedHawk 运行参数。 注意：GSR 关键字 SWITCH_MODEL_TABLE：用于设置与 SWITCH、MODEL、TABLE 相关的 RedHawk 运行参数。

\begin{noteBox}
GSR 关键字 SWITCH\_MODEL\_TABLE：用于设置与 SWITCH、MODEL、TABLE 相关的 RedHawk 运行参数。
\end{noteBox}

\section{pt2timing}
\subsection*{pt2timing}

PrimeTime TCL 接口：生成各实例的转换时间与 timing window。

GSR 关键字 ：用于设置与  相关的 RedHawk 运行参数。

GSR 关键字 ：用于设置与  相关的 RedHawk 运行参数。

\begin{lstlisting}
                       VP_PAIRS {
                       <vdd1_int> <vdd1_ext>             <vdd2_int> <vdd2_ext>
                       }
\end{lstlisting}

GSR 关键字 ：用于设置与  相关的 RedHawk 运行参数。

GSR 关键字 ：用于设置与  相关的 RedHawk 运行参数。

GSR 关键字 ：用于设置与  相关的 RedHawk 运行参数。

\begin{lstlisting}
             • the LEF file, <top_cell>_adsgds.lef, contains the cell abstractions (such as sense
\end{lstlisting}

GSR 关键字 ：用于设置与  相关的 RedHawk 运行参数。

\begin{lstlisting}
             • the DEF file, <top_cell>_adsgds.def, contains the placement of the cell and the
\end{lstlisting}

GSR 关键字 ：用于设置与  相关的 RedHawk 运行参数。

\begin{lstlisting}
             • the Synopsys LIB file, <top_cell>_adsgds.pratio, contains the relative power
\end{lstlisting}

GSR 关键字 ：用于设置与  相关的 RedHawk 运行参数。

GSR 关键字 ：用于设置与  相关的 RedHawk 运行参数。

GSR 关键字 ：用于设置与  相关的 RedHawk 运行参数。

GSR 关键字 ：用于设置与  相关的 RedHawk 运行参数。

\begin{lstlisting}
                 [ gds2rh | gds2def] -m <memory_config_file> ?<GDSII_filename>?
                 -log <log_dir_name> -out_dir <dir_name>
\end{lstlisting}

GSR 关键字 ：用于设置与  相关的 RedHawk 运行参数。

\begin{lstlisting}
                      <memory_config_file> : specifies the memory configuration file
                      <GDSII_filename> : specifies the GDS filename
                      -log <log_dir_name>: specifies the directory into which all log files are written
                      -out_dir <dir_name>: specifies the directory into which all output files are written.
\end{lstlisting}

GSR 关键字 ：用于设置与  相关的 RedHawk 运行参数。

GSR 关键字 ：用于设置与  相关的 RedHawk 运行参数。

GSR 关键字 ：用于设置与  相关的 RedHawk 运行参数。

\section{sim2iprof}
\subsection*{sim2iprof}

从仿真输出生成电流轮廓（read/write/standby 波形）。

GSR 关键字 ：用于设置与  相关的 RedHawk 运行参数。

GSR 关键字 ：用于设置与  相关的 RedHawk 运行参数。

\begin{lstlisting}
                     <design_name>.timing
\end{lstlisting}

GSR 关键字 ：用于设置与  相关的 RedHawk 运行参数。

GSR 关键字 ：用于设置与  相关的 RedHawk 运行参数。

GSR 关键字 ：用于设置与  相关的 RedHawk 运行参数。

\begin{lstlisting}
            memories, and generates current profiles in a <cell>.current file for RedHawk power
\end{lstlisting}

GSR 关键字 ：用于设置与  相关的 RedHawk 运行参数。

\begin{lstlisting}
            The Vdd current waveform is extracted and converted to RedHawk <cell>.current
\end{lstlisting}

GSR 关键字 ：用于设置与  相关的 RedHawk 运行参数。

\begin{lstlisting}
            leakage power data, and write the data to the <cell>.cdev file. Alternatively, if you have
\end{lstlisting}

GSR 关键字 ：用于设置与  相关的 RedHawk 运行参数。

\begin{lstlisting}
            the information in the config file, and sim2iprof then writes them into the <cell>.cdev file.
\end{lstlisting}

GSR 关键字 ：用于设置与  相关的 RedHawk 运行参数。

GSR 关键字 ：用于设置与  相关的 RedHawk 运行参数。

\begin{lstlisting}
                 sim2iprof <config_file> ? -log <log_file> ?
                      ?-o <current_output_file>? ?-avm <avm_config>?
\end{lstlisting}

GSR 关键字 ：用于设置与  相关的 RedHawk 运行参数。

\begin{lstlisting}
                    <config_file> : specifies the name of the configuration file
                    -log <log_file> : allows specifying a log file for each SIM2IPROF run
                    -o <current_output_file> : specifies optional output filename. The default name is
                        <cell>.current when the keyword SIM_FILE is used in the config file, and
                        <cellname>.spiprof when the keyword CELL is used.
                    -avm <avm_config> : specifies optional AVM configuration filename to allow AVM
\end{lstlisting}

GSR 关键字 ：用于设置与  相关的 RedHawk 运行参数。（默认：is the log file, sim2iprof）

\section{aplreader}
\subsection*{aplreader}

将电流轮廓数据转换为标准 APL 输出格式。

GSR 关键字 ：用于设置与  相关的 RedHawk 运行参数。

GSR 关键字 ：用于设置与  相关的 RedHawk 运行参数。

\begin{lstlisting}
            The aplreader utility extracts data from the APL waveform files such as <cell>.current,
\end{lstlisting}

GSR 关键字 ：用于设置与  相关的 RedHawk 运行参数。

GSR 关键字 ：用于设置与  相关的 RedHawk 运行参数。

GSR 关键字 ：用于设置与  相关的 RedHawk 运行参数。

\begin{lstlisting}
                 aplreader <input_dir/file> [-l <cell_list_file>]
                      -timing <apl_file_name>
\end{lstlisting}

GSR 关键字 ：用于设置与  相关的 RedHawk 运行参数。

\begin{lstlisting}
                    <input_dir/file> : specifies the input file directory and filename
                    -l <cell_list_file> : specifies an optional cell list file for which only cells in the file
\end{lstlisting}

GSR 关键字 ：用于设置与  相关的 RedHawk 运行参数。

\begin{lstlisting}
                    -timing : prints timing info such as delay/slew/firing in the specified file
\end{lstlisting}

GSR 关键字 ：用于设置与  相关的 RedHawk 运行参数。

GSR 关键字 ：用于设置与  相关的 RedHawk 运行参数。

\begin{lstlisting}
            -----------------------------------------
            cell : <cellname> VG | LG
            vdd1 = <vdd value 1> ; vdd2 = <vdd value 2>; ... vddN = <vdd value
            N) ; C1 = <Cload1 value> ; R = <load resistance> ; C2 = <Cload2
            value>; Slew1 = <rising transition time>; Slew2 = <falling
            transition time> ;
            state = <state name> ; vector = <Boolean_eq> ; active_input =
            <active_input name> ; active_output = <active_output name> ;
\end{lstlisting}

GSR 关键字 ：用于设置与  相关的 RedHawk 运行参数。

\begin{lstlisting}
            <pin name> <peak I - A> <Total charge - C> <Wave duration - S>
\end{lstlisting}

\section{aplcdev2pwc}
\subsection*{aplcdev2pwc}

校验 cdev 数据有效性并转换为 pwcdev 格式。

GSR 关键字 ：用于设置与  相关的 RedHawk 运行参数。

GSR 关键字 ：用于设置与  相关的 RedHawk 运行参数。（默认：_STATE_HIGH）

\begin{lstlisting}
            -----------------------------------
\end{lstlisting}

GSR 关键字 ：用于设置与  相关的 RedHawk 运行参数。

GSR 关键字 ：用于设置与  相关的 RedHawk 运行参数。

\begin{lstlisting}
                     aplcdev2pwc [-cf <cap_factor>] [-rf <res_factor>] [-lf
                     <leakage_factor>] [-o <output_pwcdev_file>]
                     <input_APL_config_file> <input_cdev_file>
\end{lstlisting}

GSR 关键字 ：用于设置与  相关的 RedHawk 运行参数。

\begin{lstlisting}
                        -cf: specifies the capacitance factor (default 0.5)
                        -rf: specifies the resistance factor (default -0.5)
                        -lf: specifies the leakage current factor (default 1.5)
                        -o: specifies the output pwcdev filename
\end{lstlisting}

GSR 关键字 ：用于设置与  相关的 RedHawk 运行参数。

GSR 关键字 ：用于设置与  相关的 RedHawk 运行参数。

GSR 关键字 ：用于设置与  相关的 RedHawk 运行参数。

\section{aplcopy}
\subsection*{aplcopy}

复制并编辑 *spiprof、*cdev、*pwcdev 文件中的数据。

GSR 关键字 ：用于设置与  相关的 RedHawk 运行参数。

\begin{lstlisting}
                    PROCESS <process name> - specifies a process type, such as TT, SS, or FF
                    VDD <value> - specifies the voltage
                    VDD_PIN_NAME <name> - specifies the VDD pin name
                    GND_PIN_NAME <name> - specifies the ground pin name
                    PRIMARY_VDD_PIN <name> - in low power designs, specifies the associated
\end{lstlisting}

GSR 关键字 ：用于设置与  相关的 RedHawk 运行参数。

\begin{lstlisting}
                    PRIMARY_GND_PIN <name> - in low power designs, specifies the associated
\end{lstlisting}

GSR 关键字 ：用于设置与  相关的 RedHawk 运行参数。

\begin{lstlisting}
                    SWEEPVALUE <value1> [<value2> ...] - specifies a list of voltage samples in
\end{lstlisting}

GSR 关键字 ：用于设置与  相关的 RedHawk 运行参数。

\begin{lstlisting}
                    LEAKAGE_SCALING <leakage_scaling_factor_for_vdd1>
                       [<leakage_scaling_factor_for_vdd2> ... ] - specifies custom scaling factors
\end{lstlisting}

GSR 关键字 ：用于设置与  相关的 RedHawk 运行参数。

GSR 关键字 ：用于设置与  相关的 RedHawk 运行参数。

GSR 关键字 ：用于设置与  相关的 RedHawk 运行参数。

\begin{lstlisting}
            *spiprof, *cdev and *pwcdev files using a mapping (configuration) file.
\end{lstlisting}

GSR 关键字 ：用于设置与  相关的 RedHawk 运行参数。

GSR 关键字 ：用于设置与  相关的 RedHawk 运行参数。

\begin{lstlisting}
                     aplcopy -c -pwc <input_file/directory> <output_file>
                     ?–ilimit? ?[-l <list_file>]? ?[-map <map_file>]?
                     <input.current> <output.current>
\end{lstlisting}

GSR 关键字 ：用于设置与  相关的 RedHawk 运行参数。

\begin{lstlisting}
                    -c: specifies the source is a cdev file/directory.
                    -pwc: specifies the source is a pwcdev file/directory.
                    -input_file/directory: specifies the source filename or directory name to be edited
\end{lstlisting}

GSR 关键字 ：用于设置与  相关的 RedHawk 运行参数。

\begin{lstlisting}
                    -output_file: specifies the output filename.
                    -ilimit: specifies no limit data checking is to be performed on the results. Optional.
                    -l <list_file>: specifies only cells in the <list_file> are to be edited or copied to
                    the <output_file>. Without this option, all input cells are copied to the output
\end{lstlisting}

GSR 关键字 ：用于设置与  相关的 RedHawk 运行参数。（可选）

\begin{lstlisting}
                    -map <map_file>: specifies the configuration map file containing editing details for
\end{lstlisting}

GSR 关键字 ：用于设置与  相关的 RedHawk 运行参数。

\begin{lstlisting}
                    does not edit the source file/directory, or copies selected cells if a <list_file>
\end{lstlisting}

GSR 关键字 ：用于设置与  相关的 RedHawk 运行参数。（可选）

\section{aplchk}
\subsection*{aplchk}

对 APL 电容与 PWC 数据进行有效性检查。

GSR 关键字 ：用于设置与  相关的 RedHawk 运行参数。

GSR 关键字 ：用于设置与  相关的 RedHawk 运行参数。（可选）

\begin{lstlisting}
                      TEMP <new temp value> #optional
                      MODEL <new model> #optional
                      CELL_MAPPING {
                       <original_cell_name> <new_cell_name>
                      }
                      PIN_MAPPING {
                      #atleast one line must be specified in PIN_MAPPING
                        <original pin-name1> <new pin name>
                      }
                      CDEV <pin_name> {
                      # optional. The specified pin name should be the new pin
                      name if it has been modified in PIN_MAPPING {}.
\end{lstlisting}

\begin{lstlisting}
                      RES <new_esr_value> #optional
                      CAP <new_cdev_value> #optional
                      LEAK <new_leak_value> #optional
                      }
\end{lstlisting}

GSR 关键字 ：用于设置与  相关的 RedHawk 运行参数。

GSR 关键字 ：用于设置与  相关的 RedHawk 运行参数。

GSR 关键字 ：用于设置与  相关的 RedHawk 运行参数。

\begin{lstlisting}
                    aplchk -pwc -l <cell_list> -spice <apl_config file>
\end{lstlisting}

GSR 关键字 ：用于设置与  相关的 RedHawk 运行参数。

\begin{lstlisting}
            CELL <cellname> <power/ground pin> <state> <voltage value index>
            <model peak-to-peak current value> <cell peak-to-peak current value>
            <diff = abs(model pp. value - cell pp. value) / max(cell pp. value, model pp.
            value) * 100%>
\end{lstlisting}

\section{clampviewer}
\subsection*{clampviewer}

Linux 可执行工具，预览 ESD clamp 的 I-V 曲线数据。

GSR 关键字 ：用于设置与  相关的 RedHawk 运行参数。

GSR 关键字 ：用于设置与  相关的 RedHawk 运行参数。

\begin{lstlisting}
            show clamp_iv <ivName>'. The syntax is:
                      clampviewer <IV_FileName> [<IV_Name>] [<options>]
\end{lstlisting}

GSR 关键字 ：用于设置与  相关的 RedHawk 运行参数。（默认：all I-Vs in the file are plotted）

\section{通用注意事项}
\subsection*{General Notes}

\begin{noteBox}
vcdtrans/vcdscan/fsdbtrans 的 \texttt{-s}/\texttt{-e} 时间值若小于 0.1，RedHawk 假定单位为秒；否则为皮秒（ps）。
\end{noteBox}

\begin{tipBox}
VCD 动态 IR 流程：先用 vcdscan 获取峰值功耗周期 \texttt{<FROM>,<TO>}，
再在 GUI Dynamic 菜单中启动 VCD-based Dynamic IR-drop Analysis。
\end{tipBox}
